\section{强场量子动力学}
\label{ch:gauge-volkov-sfa}

\chapref{ch:strong-field-classical}给出了强场中的经典连续态轨道。量子理论仍从精确最小耦合与规范协变出发；舍弃母离子势后才得到 Volkov 态，再用 Volkov 传播替代精确连续态传播才得到 SFA。非周期电离振幅的渐近结构由 Keldysh 复时间处理，严格周期驱动则由 Floquet--Sambe 准能谱处理。
\subsection{最小耦合、规范协变与精确传播}

设质量为 \(m\)、电荷为 \(Q\) 的非相对论粒子处在外电磁势 \((\Phi,\mathbf A)\) 中。经典 Lagrangian 为

\[
L=\frac12m\dot{\mathbf r}^{\,2}
+Q\dot{\mathbf r}\cdot\mathbf A(\mathbf r,t)
-Q\Phi(\mathbf r,t)-V(\mathbf r).
\]

正则动量与机械动量不同：

\[
\mathbf p_{\rm can}
\equiv\frac{\partial L}{\partial\dot{\mathbf r}}
=m\dot{\mathbf r}+Q\mathbf A,
\qquad
\boldsymbol\pi\equiv m\dot{\mathbf r}
=\mathbf p_{\rm can}-Q\mathbf A.
\]

Legendre 变换给出

\[
H
=\frac{[\mathbf p_{\rm can}-Q\mathbf A(\mathbf r,t)]^2}{2m}
+Q\Phi(\mathbf r,t)+V(\mathbf r).
\]

正则量子化 \(\mathbf p_{\rm can}\mapsto\hat{\mathbf p}=-\ii\hbar\nabla\) 后得到最小耦合哈密顿量

\begin{equation}
\hat H(t)
=\frac{1}{2m}
[\hat{\mathbf p}-Q\mathbf A(\mathbf r,t)]^2
+Q\Phi(\mathbf r,t)+V(\mathbf r).
\label{eq:minimal-coupling}
\end{equation}

\eqnrefs{eq:minimal-coupling}是后续所有规范选择的共同起点。对电子 \(Q=Q_e=-e\)、质量为 \(m_e\)。只有在这一一般形式建立以后，我们才进入强场常用的偶极近似 \(k a\ll1\)，把原子尺度内的矢势写成 \(\mathbf A(\mathbf r,t)\simeq\mathbf A(t)\)。在速度规范且取 \(\Phi=0\) 时，电子哈密顿量为

\[
\hat H_V(t)
=\frac{[\hat{\mathbf p}+e\mathbf A(t)]^2}{2m_e}+V(\mathbf r).
\]

这里 \(\mathbf A\) 始终是 SI 矢势，量纲为动量/电荷；因此 \(e\mathbf A\) 与 \(\hat{\mathbf p}\) 具有相同的动量量纲。本章后续所有作用量均保留 \(\hbar\)，不切换单位制。

采用全书四维约定 $x^\mu=(ct,\mathbf r)$、$A_\mu=(\Phi/c,-\mathbf A)$。规范变换写成
\begin{equation}
A'_\mu=A_\mu-\partial_\mu\chi,
\qquad
\psi'=\exp\!\left(\frac{\ii Q\chi}{\hbar}\right)\psi,
\label{eq:strong-gauge-covariant-transform}
\end{equation}
其三维分量正是
\[
\mathbf A'=\mathbf A+\nabla\chi,
\qquad
\Phi'=\Phi-\partial_t\chi.
\]
因此四维式中的负号与三维矢势前的正号并不矛盾，而是 $A_i=-\mathbf A_i$ 的直接结果。

为了使协变动量 \(\hat{\mathbf p}-Q\mathbf A\) 的作用保持等价，波函数必须同时作

\begin{equation}
\psi'(\mathbf r,t)
=
\exp\!\left(\frac{\ii Q}{\hbar}\chi(\mathbf r,t)\right)
\psi(\mathbf r,t).
\label{eq:gauge-wavefunction}
\end{equation}

这是一个位置和时间依赖的幺正变换 \(U_\chi\)。时间依赖幺正变换下哈密顿量不是简单的 \(U\hat H U^\dagger\)，而是

\[
\hat H'
=U_\chi\hat H U_\chi^\dagger
+\ii\hbar(\partial_tU_\chi)U_\chi^\dagger.
\]

第二项正好产生标势的 \(-\partial_t\chi\) 修正。由此可直接验证 Schrödinger 方程的形式不变，而概率密度、概率流和其他规范不变量给出相同可观测结果。

在偶极近似下选择

\[
\chi(\mathbf r,t)=-\mathbf A(t)\cdot\mathbf r.
\]

对电子 \(Q=-e\)，利用 \(\mathbf E(t)=-\partial_t\mathbf A(t)\)，得到长度规范

\[
\hat H_L(t)
=\frac{\hat{\mathbf p}^{\,2}}{2m_e}
+V(\mathbf r)
+e\,\mathbf r\cdot\mathbf E(t).
\]

两种规范下的波函数关系为

\[
\psi_L
=\exp\!\left(\frac{\ii e}{\hbar}\mathbf A(t)\cdot\mathbf r\right)\psi_V,
\qquad
\psi_V
=\exp\!\left(-\frac{\ii e}{\hbar}\mathbf A(t)\cdot\mathbf r\right)\psi_L.
\]

后文若在有限基组或近似传播子中比较两个规范，需要记住：精确理论规范等价，并不保证任意截断近似在数值上严格规范不变。

\begin{note}[Gauge truncation]
精确理论中，由规范变换联系的哈密顿量彼此幺正等价。有限基组截断或不一致的 SFA 分割可能破坏这种等价性；因此，近似计算中的规范 dependence 应被视为近似误差的诊断，而不是新的可观测量。
\end{note}

精确规范协变不仅约束瞬时哈密顿量，也约束有限时间传播子。若 $U_\chi(t)=\exp[\ii Q\chi(\mathbf r,t)/\hbar]$，则
\begin{equation}
U'(t,t_0)=U_\chi(t)U(t,t_0)U_\chi^\dagger(t_0).
\label{eq:gauge-covariant-propagator}
\end{equation}
因此位置表象传播核满足
\begin{equation}
K_{A-\partial\chi}(x_f,x_i)
=\exp\!\left[\frac{\ii Q}{\hbar}(\chi_f-\chi_i)\right]K_A(x_f,x_i),
\label{eq:gauge-propagator-endpoint-phase}
\end{equation}
其中 $A-\partial\chi$ 使用四维记号。任何有限基或 SFA 截断若不同时变换态、算符与边界相位，就可能人为产生规范 dependence。

\begin{proof}[解]
从传播子定义直接计算。

态的规范变换。规范变换 \(U_\chi(t)=\exp[\ii Q\chi(\mathbf r,t)/\hbar]\) 把 Schrödinger 方程
\[
 \ii\hbar\partial_t|\psi(t)\rangle=\hat H_A(t)|\psi(t)\rangle
\]
变为
\[
 \ii\hbar\partial_t|\psi'(t)\rangle=\hat H_{A-\partial\chi}(t)|\psi'(t)\rangle,
 \qquad |\psi'(t)\rangle=U_\chi(t)|\psi(t)\rangle.
\]

传播子的协变变换。由 \(|\psi(t)\rangle=U(t,t_0)|\psi(t_0)\rangle\) 和 \(U_\chi\) 的幺正性，
\[
 \begin{aligned}
 |\psi'(t)\rangle
 &=U_\chi(t)|\psi(t)\rangle
 =U_\chi(t)U(t,t_0)|\psi(t_0)\rangle\\
 &=U_\chi(t)U(t,t_0)U_\chi^\dagger(t_0)U_\chi(t_0)|\psi(t_0)\rangle\\
 &=U_\chi(t)U(t,t_0)U_\chi^\dagger(t_0)|\psi'(t_0)\rangle.
 \end{aligned}
\]
该关系对任意初态 \(|\psi'(t_0)\rangle\) 成立，故传播子必须满足
\[
 U'(t,t_0)=U_\chi(t)U(t,t_0)U_\chi^\dagger(t_0),
\]
即\eqnrefs{eq:gauge-covariant-propagator}。验证：将上式代入变换后的 Schrödinger 方程，
\[
 \ii\hbar\partial_tU_\chi UU_\chi^\dagger
 =(\ii\hbar\dot U_\chi U_\chi^\dagger)U_\chi UU_\chi^\dagger
 +U_\chi(\hat H_A)U_\chi^\dagger\cdot U_\chi UU_\chi^\dagger
 =\hat H_{A-\partial\chi}\,U',
\]
其中 \(\ii\hbar\dot U_\chi U_\chi^\dagger=-Q\partial_t\chi\) 恰好把标势从 \(A_0\) 移到 \(A_0-\partial_t\chi\)。

位置表象的端点相位。夹上位置本征态 \(|\mathbf r_f\rangle\)、\(|\mathbf r_i\rangle\)，利用
\[
 \langle\mathbf r|U_\chi(t)
 =\ee^{\ii Q\chi(\mathbf r,t)/\hbar}\langle\mathbf r|,
 \qquad
 U_\chi^\dagger(t_0)|\mathbf r\rangle
 =\ee^{-\ii Q\chi(\mathbf r,t_0)/\hbar}|\mathbf r\rangle,
\]
得到
\[
 K_{A-\partial\chi}(\mathbf r_f,t_f;\mathbf r_i,t_i)
 =\ee^{\ii Q\chi_f/\hbar}\,K_A(\mathbf r_f,t_f;\mathbf r_i,t_i)\,\ee^{-\ii Q\chi_i/\hbar}
 =\exp\!\left[\frac{\ii Q}{\hbar}(\chi_f-\chi_i)\right]K_A(\mathbf r_f,t_f;\mathbf r_i,t_i),
\]
即\eqnrefs{eq:gauge-propagator-endpoint-phase}。端点相位不是额外动力学，而是同一物理传播在不同局域相位基中的坐标变换。
\end{proof}

规范协变传播子是精确的，但实际计算需要具体参考态。在强场中，自由电子在激光场下的精确解——Volkov 态——提供了分析电离和再碰撞的出发点。强场近似 (SFA) 用 Volkov 传播子替代精确传播子，将电离振幅化为鞍点积分。

\subsection{Volkov 连续态、SFA 与量子轨道}

忽略原子势 \(V(\mathbf r)\)，自由电子在激光场中的哈密顿量为

\[
\hat H_V(t)=\frac{[\hat{\mathbf p}+e\mathbf A(t)]^2}{2m_e}.
\]

由于偶极近似下 \(\mathbf A(t)\) 与 \(\hat{\mathbf p}\) 对易，可以取以最终漂移动量 \(\mathbf p\) 标记的平面波形式

\begin{equation}
\psi_{\mathbf p}^{V}(\mathbf r,t)
=(2\pi\hbar)^{-3/2}
\exp\!\left(\frac{\ii}{\hbar}\mathbf p\cdot\mathbf r\right)
\exp\!\left[-\frac{\ii}{\hbar}\mathcal S_V(\mathbf p,t)\right],
\label{eq:volkov-action}
\end{equation}


\par
\begin{workedexample}[圆偏振场中的 Volkov 准能]
取圆偏振矢势 $\vb A(t)=\dfrac{E_0}{\omega}\bigl(\cos\omega t\,\hat{\vb x}+\sin\omega t\,\hat{\vb y}\bigr)$，证明 Volkov 作用量\eqnrefs{eq:volkov-action}中的被积函数为常数，从而圆偏振场中电子准能恰为 $\varepsilon_p=p^2/(2m_e)+2U_p$；说明一般横向动量下 \(\\vb{p}\\cdot\\vb{A}(t)\) 仍产生周期微动。

\textbf{解.} 对圆偏振 \(\mathbf A=A_0(\cos\omega t,\sin\omega t)\)，有 \(A^2=A_0^2\) 恒定。展开为 \([\vb p+e\vb A(t)]^2=p^2+2e\vb p\cdot\vb A(t)+e^2A_0^2\)；但 \(\mathbf p\cdot\mathbf A(t)\) 对一般横向 \(\mathbf p\) 仍随时间振荡。因此作用量由平均项和周期项组成；周期项只改变 Volkov 态的周期微动，不改准能。平均动能为
\[
\overline{\frac{[\mathbf p+e\mathbf A(t)]^2}{2m_e}}=\frac{p^2}{2m_e}+\frac{e^2A_0^2}{2m_e}=\frac{p^2}{2m_e}+2U_p,
\]
其中 \(U_p=e^2E_0^2/(4m_e\omega^2)\)。故题设的准能结论成立，但“被积函数恒定、无时间振荡项”仅在 \(\mathbf p\) 沿偏振平面法向时成立；一般情形须保留周期微动项。
\end{workedexample}
\par
其中

\[
\mathcal S_V(\mathbf p,t)
=
\int^t\dd\tau\,
\frac{[\mathbf p+e\mathbf A(\tau)]^2}{2m_e}
\]

称为 Volkov 作用量（Volkov 作用量），量纲为作用量。

验证时保留 \(\hbar\)：

\[
\ii\hbar\frac{\partial\psi}{\partial t}
=\dot{\mathcal S}_V\psi
=\frac{[\mathbf p+e\mathbf A(t)]^2}{2m_e}\psi.
\]

设

\[
\mathbf A(t)=-\frac{\mathbf E_0}{\omega}\sin\omega t.
\]

则

\[
\frac{[\mathbf p+e\mathbf A(t)]^2}{2m_e}
=\frac{p^2}{2m_e}
-\frac{e\,\mathbf p\cdot\mathbf E_0}{m_e\omega}\sin\omega t
+\frac{e^2E_0^2}{2m_e\omega^2}\sin^2\omega t.
\]

利用

\[
\sin^2\omega t=\frac12(1-\cos2\omega t),
\]

周期平均项为

\[
\frac{p^2}{2m_e}+\frac{e^2E_0^2}{4m_e\omega^2}
=\frac{p^2}{2m_e}+U_p.
\]

因此 Volkov 态带有平均准能量

\[
\varepsilon_p=\frac{p^2}{2m_e}+U_p.
\]

其中
\begin{equation}
U_p=\frac{e^2E_0^2}{4m_e\omega^2}
\label{eq:ponderomotive-energy-restated}
\end{equation}
是有质动力能——电子在激光场中抖动运动的周期平均动能。

\begin{proof}[解]
经典电子在平面波场 \(\mathbf E(t)=\mathbf E_0\cos\omega t\) 中的运动方程为 \(m_e\ddot{\mathbf x}=-e\mathbf E_0\cos\omega t\)。对时间积分两次（设初速度为零、初位置为零）：
\begin{equation*}
\dot{\mathbf x}(t)=-\frac{e\mathbf E_0}{m_e\omega}\sin\omega t,\qquad
\mathbf x(t)=\frac{e\mathbf E_0}{m_e\omega^2}\sin\omega t\equiv\boldsymbol\alpha_0\sin\omega t,
\end{equation*}
其中 \(\alpha_0=eE_0/(m_e\omega^2)\) 为抖动振幅。

抖动动能为 \(\frac{1}{2}m_e\dot{\mathbf x}^2=\frac{1}{2}m_e\alpha_0^2\omega^2\cos^2\omega t\)。周期平均利用 \(\langle\cos^2\omega t\rangle=\frac{1}{2}\)：
\begin{equation*}
\frac12 m_e\langle\dot{\mathbf x}^2\rangle
=\frac12 m_e\alpha_0^2\omega^2\cdot\frac12
=\frac{m_e\alpha_0^2\omega^2}{4}
=\frac{m_e}{4}\cdot\frac{e^2E_0^2}{m_e^2\omega^4}\cdot\omega^2
=\frac{e^2E_0^2}{4m_e\omega^2}
=U_p.
\end{equation*}
即有质动力能\eqnrefs{eq:ponderomotive-energy-restated}。

电离后连续态总能量为 \(\varepsilon_p=p^2/(2m_e)+U_p\)——抖动动能 \(U_p\) 作为有效势能出现在连续态中。能量守恒要求 \(n\hbar\omega=I_p+p^2/(2m_e)+U_p\)（吸收 \(n\) 个光子，克服电离能 \(I_p\)，剩余为光电子动能加抖动能）。因此光电子动能
\begin{equation*}
\frac{p^2}{2m_e}=n\hbar\omega-I_p-U_p.
\end{equation*}
与无场情形相比阈值移动 \(-U_p\)，最小吸收光子数增至 \(n_{\min}=\lceil(I_p+U_p)/\hbar\omega\rceil\)。

Keldysh 参数 \(\gamma_K=\sqrt{I_p/(2U_p)}\) 划分两区：\(\gamma_K\gg1\) 为多光子区（微扰论适用），\(\gamma_K\ll1\) 为隧穿区（ADK 适用）。

抖动能的空间梯度给出 \(\mathbf F_p=-\nabla U_p\)，电子被推出高光强区域。

\end{proof}

这解释了 ATI 光电子峰中为什么出现 \(-U_p\) 的阈值移动。

Volkov 态给出了自由电子在激光场中的精确解。将精确传播子分解为 Volkov 传播子加 Coulomb 修正，在核外区域忽略 Coulomb 势即得强场近似 (SFA)。

完整哈密顿量仍写成

\[
H(t)=H_0+H_I(t),
\qquad
H_0=\frac{\hat{\mathbf p}^{\,2}}{2m_e}+V(\mathbf r).
\]

相互作用表象的传播子具有形式精确的 Dyson 表达，这一点已在\chapref{ch:time-dependent-perturbation}中建立。强场问题的困难在于：电离前原子势不能忽略，而电子进入连续谱以后激光场又不能作为小微扰处理。因此 SFA 不是简单的“对场做低阶展开”，而是重新选择零阶传播：束缚阶段保留原子本征态，连续阶段改用 Volkov 传播子。

设初态在脉冲到来前为

\[
|\psi_0(t)\rangle=\ee^{+\ii I_p t/\hbar}|\psi_0\rangle,
\qquad H_0|\psi_0\rangle=-I_p|\psi_0\rangle.
\]

在长度规范中，直接电离的 SFA 振幅写成

\begin{equation}
M_{\rm SFA}(\mathbf p)
=-\frac{\ii}{\hbar}\int_{t_0}^{T}\dd t'
\,\langle\chi_{\mathbf p}^{V}(t')|
H_I(t')|\psi_0(t')\rangle,
\label{eq:sfa-direct-amplitude}
\end{equation}

其中 \(|\chi_{\mathbf p}^{V}\rangle\) 是与最终漂移动量 \(\mathbf p\) 对应的 Volkov 态。\eqnrefs{eq:sfa-direct-amplitude}包含三项物理近似：初态 depletion 被忽略；电子离开束缚区后用 Volkov 态代替真实 Coulomb 连续态；电子在连续传播期间不再被母离子势显著偏转。最后一项正是普通 SFA 在低能电子和长程 Coulomb 聚焦问题上产生系统误差的来源。

把初态相位和 Volkov 相位合并，振幅可整理为

\[
M_{\rm SFA}(\mathbf p)
=\int_{t_0}^{T}\dd t'\,
D(\mathbf p,t')
\exp\!\left[\frac{\ii}{\hbar}\mathcal S(\mathbf p,t';T)\right],
\]

其中

\begin{equation}
\mathcal S(\mathbf p,t';T)
=-\int_{t'}^{T}\dd\tau
\left[
\frac{[\mathbf p+e\mathbf A(\tau)]^2}{2m_e}+I_p
\right]
\label{eq:sfa-action}
\end{equation}

而 \(D(\mathbf p,t')\) 包含偶极矩阵元、场包络以及可能的角动量因子。对有限脉冲，\(T\) 取在脉冲结束以后；对严格无限长周期场，则应转到 Floquet 准能量语言，而不能把一个线性随时间发散的 Volkov 相位直接解释为普通散射相位。

SFA 没有一个在所有强场参数区间都普适的小参数。它通常在电子离核后抖动动能远大于 Coulomb 势修正、且目标可观测量不强烈依赖低能长程散射时最可靠。若需要系统恢复 Coulomb 效应，可以进一步使用 Coulomb--Volkov、eikonal-Volkov 或 Coulomb-corrected SFA 等方法。

当作用量在一个光学周期内改变很多个弧度时，\eqnrefs{eq:sfa-direct-amplitude}的主要贡献来自复时间平面中的驻相点。由

\[
\frac{\partial \mathcal S}{\partial t_s}=0
\]

得到

\begin{equation}
\frac{[\mathbf p+e\mathbf A(t_s)]^2}{2m_e}+I_p=0.
\label{eq:sfa-saddle}
\end{equation}

若 \(t_s\) 为实数，则第一项非负，而 \(I_p>0\)，因此左边不可能为零。故束缚态电离的主鞍点必然满足

\[
t_s\in\mathbb C.
\]

这给出了复时间的直接动力学来源：电子从束缚态到达经典允许连续谱之前，必须经过一个没有实经典轨道对应的禁阻段。

仅求出鞍点位置还不够。把作用量在每一个非简并鞍点附近展开，

\[
\mathcal S(t)
=\mathcal S(t_s)
+\frac12\mathcal S''(t_s)(t-t_s)^2+\cdots,
\]

其中

\[
\mathcal S''(t_s)
=\frac{e}{m_e}[\mathbf p+e\mathbf A(t_s)]\cdot\dot{\mathbf A}(t_s)
=-\frac{e}{m_e}[\mathbf p+e\mathbf A(t_s)]\cdot\mathbf E(t_s).
\]

沿最陡下降路径作高斯积分，单个鞍点的贡献具有结构

\begin{equation}
M_s(\mathbf p)
\simeq
D(\mathbf p,t_s)
\left[
\frac{2\pi\ii\hbar}{\mathcal S''(t_s)}
\right]^{1/2}
 \ee^{\ii\mathcal S(\mathbf p,t_s)/\hbar}.
\label{eq:sfa-saddle-prefactor}
\end{equation}

平方根的相位由实际最陡下降轮廓决定。一个激光周期通常存在多个相关鞍点，因此完整振幅是

\[
M(\mathbf p)\simeq\sum_s M_s(\mathbf p).
\]

不同复鞍点之间的相位差会产生 ATI 峰内结构、半周期干涉以及所谓 intracycle/intercycle interference。所谓“量子轨道”并不是电子具有一条可直接观测的经典路径，而是每一个复鞍点及其对应的半经典作用量对积分作出的驻相贡献。

若讨论 HHG 或 rescattering ATI，还必须在直接电离振幅之外加入电子从电离时刻 \(t'\) 传播到回碰时刻 \(t\) 的连续态传播。中间漂移动量 \(\mathbf p\) 也要积分。典型半经典作用量为

\[
\mathcal S(\mathbf p,t,t')
=
-\int_{t'}^{t}\dd\tau
\left[
\frac{[\mathbf p+e\mathbf A(\tau)]^2}{2m_e}+I_p
\right]
+\hbar\Omega t,
\]

其中 \(\hbar\Omega\) 是最终发射光子的能量。对 \(t'\)、\(\mathbf p\) 与 \(t\) 同时驻相，分别得到

\begin{align*}
\frac{[\mathbf p_s+e\mathbf A(t_s')]^2}{2m_e}+I_p&=0,
\\
\int_{t_s'}^{t_s}\dd\tau\,\frac{\mathbf p_s+e\mathbf A(\tau)}{m_e}&=0,
\\
\hbar\Omega&=I_p+\frac{[\mathbf p_s+e\mathbf A(t_s)]^2}{2m_e}.
\end{align*}

第一式是复时间电离条件；第二式是电子位移回到母离子位置的条件；第三式是复合时的能量守恒。把 \(I_p\to0\) 并把时间取成实数，就退化为\chapref{ch:strong-field-classical}的经典三步模型。因此经典的 \(3.17U_p\) 截止并不是与 SFA 无关的额外规则，而是这一组量子鞍点方程的实轨道极限。

\begin{note}[复时间与“隧穿时间”]
鞍点的虚部控制指数抑制，但不能在没有额外测量定义时直接等同于唯一的可观测“隧穿时间”。不同操作性定义可能对应 Wigner time、Larmor time、attoclock offset 等不同量。SFA 的复时间首先是解析延拓后的驻相坐标。
\end{note}

SFA 的鞍点方程已经给出复电离时间，但不同渐近区间的物理图像差异需要通过 Keldysh 参数 $\gamma_K$ 统一描述。以下从鞍点方程出发，导出有限 $\gamma_K$ 的复时间作用量，再分别取隧穿和多光子极限。

\subsection{Keldysh 复时间渐近}\label{sec:keldysh-complex-time}

本章前文的 SFA 鞍点方程已经表明，束缚态电子的电离时刻一般不是实数。Keldysh 理论的核心不是简单地把强场电离分成“隧穿”和“多光子”两类，而是把这两种极限放在同一个复时间作用量中，并找出控制两种极限之间连续过渡的无量纲参数。把有限 \(\gamma_K\) 的作用量积分出来以后，\(\gamma_K\ll1\) 与 \(\gamma_K\gg1\) 才分别显出隧穿和多光子的渐近形式。

本章继续使用 SI 制。下面所有鞍点方程、作用量与渐近展开都显式保留 \(m_e,e,\hbar\)。

取线偏振单色场

\[
E(t)=E_0\cos\omega t,
\qquad
A(t)=-\frac{E_0}{\omega}\sin\omega t,
\]

并考虑最简单的纵向最终漂移动量 \(p=0\)。SFA 鞍点方程为

\[
\frac{e^2A^2(t_s)}{2m_e}+I_p=0,
\]

因此

\[
eA(t_s)=\pm \ii\sqrt{2m_eI_p}.
\]

场峰取在 \(t=0\)。主导鞍点位于上半复平面，可写成

\[
t_s=\ii\tau_0,
\qquad \tau_0>0.
\]

利用

\[
\sin(\ii\omega\tau_0)
=\ii\sinh(\omega\tau_0),
\]

得到

\[
\frac{eE_0}{\omega}\sinh(\omega\tau_0)
=\sqrt{2m_eI_p}.
\]

于是自然定义

\begin{equation}
\gamma_K
\equiv
\frac{\omega\sqrt{2m_eI_p}}{eE_0}
=\sqrt{\frac{I_p}{2U_p}},
\qquad
U_p=\frac{e^2E_0^2}{4m_e\omega^2}.
\label{eq:keldysh}
\end{equation}

鞍点的虚部因此不是独立参数，而是

\begin{equation}
\omega\tau_0=\operatorname{arsinh}\gamma_K.
\label{eq:keldysh-imag-time}
\end{equation}

当 \(\gamma_K\ll1\) 时，\(\omega\tau_0\simeq\gamma_K\ll1\)，电子穿越禁阻区的复时间短于一个光学周期，外场在这一阶段近似静止；当 \(\gamma_K\gg1\) 时，\(\tau_0\) 不再对应准静态穿隧过程，周期驱动的多光子结构变得本质。

Keldysh 参数可以直接由复鞍点的虚时间尺度理解。定义束缚能对应的隧穿动量尺度

\[
p_I\equiv\sqrt{2m_eI_p},
\]

其量纲为动量。对准静态场，主鞍点满足 \(\omega\tau_0\ll1\)，于是 \(\sinh(\omega\tau_0)\simeq\omega\tau_0\)，从前一节的精确鞍点方程得到

\[
\tau_T\equiv\tau_0\simeq\frac{p_I}{eE_0}.
\]

这个时间尺度不能简单地取成``势垒宽度除以一个固定速度''：三角势垒中的虚速度沿禁阻区变化，直接使用 \(I_p/(eE_0)\) 除以 \(p_I/m_e\) 会少一个因子 2。复时间鞍点给出的 \(p_I/(eE_0)\) 才与 WKB 作用量一致。

激光场变化时间为 \(\tau_L=\omega^{-1}\)，因此

\[
\frac{\tau_T}{\tau_L}
=\frac{\omega p_I}{eE_0}
=\gamma_K.
\]

同一个参数也可写成能量尺度之比：

\[
\gamma_K^2=\frac{I_p}{2U_p}.
\]

因此 \(\gamma_K\) 同时编码束缚能、抖动动能和驱动时间尺度。提高场强 \(E_0\) 会减小 \(\gamma_K\)，提高角频率 \(\omega\) 会增大 \(\gamma_K\)，不同原子的 \(I_p\) 也直接改变动力学区间。

仅知道 \(\gamma_K\) 的定义还不能给出电离率。真正决定指数抑制的是从复鞍点 \(t_s=\ii\tau_0\) 到实时间出口 \(t=0\) 的作用量虚部。对 \(p=0\)，把 SFA 指数中的动力学部分写成

\[
\mathcal S(t_s,0)
=\int_{t_s}^{0}
\left[
\frac{e^2A^2(t)}{2m_e}+I_p
\right]\dd t.
\]

令 \(t=\ii\tau\)，则 \(\dd t=\ii\dd\tau\)，并且

\[
A(\ii\tau)
=-\ii\frac{E_0}{\omega}\sinh\omega\tau.
\]

因此禁阻段的正实 Euclidean 作用量可以定义为

\begin{equation}
S_I
\equiv
\int_{0}^{\tau_0}
\left[
I_p-
\frac{e^2E_0^2}{2m_e\omega^2}
\sinh^2(\omega\tau)
\right]\dd\tau.
\label{eq:keldysh-euclidean-action}
\end{equation}

积分并不需要近似。利用

\[
\int\sinh^2(\omega\tau)\dd\tau
=
\frac{\sinh(2\omega\tau)}{4\omega}-\frac{\tau}{2},
\]

以及

\[
\omega\tau_0=\operatorname{arsinh}\gamma_K,
\qquad
\sinh(2\omega\tau_0)
=2\gamma_K\sqrt{1+\gamma_K^2},
\]

再代入 \(U_p=e^2E_0^2/(4m_e\omega^2)=I_p/(2\gamma_K^2)\)，得到

\begin{equation}
S_I
=
\frac{I_p}{\omega}
\left[
\left(1+\frac{1}{2\gamma_K^2}\right)
\operatorname{arsinh}\gamma_K
-
\frac{\sqrt{1+\gamma_K^2}}{2\gamma_K}
\right].
\label{eq:keldysh-action-finite-gamma}
\end{equation}

在鞍点近似中，电离振幅的指数部分为 \(\exp(-S_I/\hbar)\)，因此概率或速率的主指数为

\begin{equation}
W_{\rm exp}
\propto
\exp\left\{
-\frac{2I_p}{\hbar\omega}
 f(\gamma_K)
\right\},
\quad
f(\gamma_K)
=
\left(1+\frac{1}{2\gamma_K^2}\right)
\operatorname{arsinh}\gamma_K
-
\frac{\sqrt{1+\gamma_K^2}}{2\gamma_K}.
\label{eq:keldysh-exponent}
\end{equation}

这里写的是指数精度。真实电离率还包含由束缚态结构、角动量、Coulomb 尾以及鞍点二阶涨落决定的预因子。Keldysh 原始理论、PPT 理论与 ADK 理论之间的重要差异，主要就体现在如何处理这些预因子以及 Coulomb 修正。

\begin{proof}[解]
换元到无量纲变量。从 \eqnrefs{eq:keldysh-euclidean-action} 出发，令 \(u=\omega\tau\)、\(u_0=\operatorname{arsinh}\gamma_K\)（由\eqnrefs{eq:keldysh-imag-time}给出）。则
\[
S_I
=\frac{I_p}{\omega}u_0
-
\frac{e^2E_0^2}{2m_e\omega^3}
\int_0^{u_0}\sinh^2u\,\dd u.
\]
双曲积分。利用恒等式 $\sinh^2u=(\cosh 2u-1)/2$，
\[
\int_0^{u_0}\sinh^2u\,\dd u
=\frac14\sinh 2u_0-\frac12u_0.
\]
再用 $\sinh 2u_0=2\sinh u_0\cosh u_0=2\gamma_K\sqrt{1+\gamma_K^2}$，代入得
\[
S_I
=\frac1\omega
\left[
I_pu_0+U_pu_0-\frac{U_p}{2}\sinh 2u_0
\right],
\]
用 \(\gamma_K\) 表达。代入
\[
U_p=\frac{I_p}{2\gamma_K^2},
\qquad
\sinh u_0=\gamma_K,
\qquad
\cosh u_0=\sqrt{1+\gamma_K^2},
\]
把 $u_0=\operatorname{arsinh}\gamma_K$ 和 $\sinh 2u_0=2\gamma_K\sqrt{1+\gamma_K^2}$ 代入，整理得
\[
S_I
=\frac{I_p}{\omega}
\left[
\left(1+\frac{1}{2\gamma_K^2}\right)\operatorname{arsinh}\gamma_K
-\frac{\sqrt{1+\gamma_K^2}}{2\gamma_K}
\right],
\]
即 \eqnrefs{eq:keldysh-action-finite-gamma}。

两个极限的渐近核对。当 $\gamma_K\ll1$，$\operatorname{arsinh}\gamma_K\approx\gamma_K-\gamma_K^3/6$，$\sqrt{1+\gamma_K^2}\approx1+\gamma_K^2/2$，代入得 $f(\gamma_K)\approx 2\gamma_K/3$，于是 $S_I\approx \frac{2I_p\gamma_K}{3\omega}=\frac{2(2m_eI_p)^{3/2}}{3m_e e E_0}$，正是静态隧穿 WKB 作用量。当 $\gamma_K\gg1$，$\operatorname{arsinh}\gamma_K\approx\ln(2\gamma_K)$，$f(\gamma_K)\approx\ln(2\gamma_K)-1/2$，$S_I\approx \frac{I_p}{\omega}(\ln(2\gamma_K)-1/2)$，电离率 $\propto\exp(-2S_I/\hbar)$ 呈幂律 $\gamma_K^{-2I_p/(\hbar\omega)}$，对应多光子区 $n=I_p/(\hbar\omega)$ 光子吸收。


\end{proof}

有了有限 $\gamma_K$ 的统一作用量，隧穿极限 $\gamma_K\ll1$ 和多光子极限 $\gamma_K\gg1$ 分别来自对 $f(\gamma_K)$ 的渐近展开。

对 \(f(\gamma_K)\) 在 \(\gamma_K=0\) 附近展开：

\[
\operatorname{arsinh}\gamma_K
=\gamma_K-\frac{\gamma_K^3}{6}+\frac{3\gamma_K^5}{40}+\cdots,
\]

\[
\sqrt{1+\gamma_K^2}
=1+\frac{\gamma_K^2}{2}-\frac{\gamma_K^4}{8}+\cdots.
\]

代回后

\[
f(\gamma_K)
=\frac{2\gamma_K}{3}
-\frac{\gamma_K^3}{15}
+\frac{3\gamma_K^5}{140}
+O(\gamma_K^7).
\]

保留最低阶并使用 \(\gamma_K=\omega p_I/(eE_0)\)、\(p_I=\sqrt{2m_eI_p}\)，得到

\begin{equation}
W_{\rm exp}
\propto
\exp\!\left[-\frac{2p_I^3}{3m_e e\hbar E_0}\right]
=
\exp\!\left[-\frac{2(2m_eI_p)^{3/2}}{3m_e e\hbar E_0}\right].
\label{eq:static-tunnel-exponent}
\end{equation}

这正是准静态三角势垒的 WKB 概率指数。换言之，静态隧穿公式不是与 Keldysh 理论并列的经验规律，而是有限频率复时间作用量在 \(\gamma_K\to0\) 的渐近极限。

\begin{exbox}[$f(\gamma_K)$ 两极限的数值检查]
以\ref{sec:strong-field-experiments}表中参数为例，直接代入\eqnrefs{eq:keldysh-exponent}。H 在 800\,nm、\(10^{15}\,\mathrm{W/cm^2}\)（\(\gamma_K=0.34\)）：精确值 \(f(0.34)=0.224\)，隧穿渐近 \(2\gamma_K/3=0.227\)，偏差仅 \(1\%\)；同一原子在 \(10^{14}\,\mathrm{W/cm^2}\)（\(\gamma_K=1.07\)）：\(f(1.07)=0.651\)，而 \(2\gamma_K/3=0.713\) 偏差已达 \(9\%\)，多光子渐近 \(\ln(2\gamma_K)-1/2=0.26\) 则完全失效。反向地，Xe 在 800\,nm、\(10^{12}\,\mathrm{W/cm^2}\)（\(\gamma_K\approx10\)）：\(f(10)=2.51\) 与 \(\ln(2\gamma_K)-1/2=2.50\) 吻合到 \(0.6\%\)。可见隧穿渐近在 \(\gamma_K\lesssim0.5\) 才定量可靠，多光子渐近则要求 \(\gamma_K\gtrsim5\)；中间区域（多数实验恰在其中）必须直接使用完整的\eqnrefs{eq:keldysh-exponent}。
\end{exbox}

为了直接核对，可在静态场近似中写

\[
|p(x)|=\sqrt{2m_e(I_p-eE_0x)},
\qquad
0<x<\frac{I_p}{eE_0},
\]

于是

\[
\frac{2}{\hbar}\int_0^{I_p/(eE_0)}|p(x)|\,\dd x
=
\frac{2(2m_eI_p)^{3/2}}{3m_e e\hbar E_0},
\]

与 \eqnrefs{eq:static-tunnel-exponent} 完全一致。前面的因子 2 来自概率为振幅模平方，因此不能把振幅指数和概率指数混在一起。

\begin{example}[多光子渐近与离散阈值]

大 \(\gamma_K\) 时

\[
\operatorname{arsinh}\gamma_K
=\ln(2\gamma_K)+O(\gamma_K^{-2}),
\]

\[
\frac{\sqrt{1+\gamma_K^2}}{2\gamma_K}
=\frac12+O(\gamma_K^{-2}),
\]

因此

\[
f(\gamma_K)
=\ln(2\gamma_K)-\frac12
+O\!\left(\frac{\ln\gamma_K}{\gamma_K^2}\right).
\]

于是指数部分

\[
W_{\rm exp}
\sim
\exp\left[
-\frac{2I_p}{\hbar\omega}
\left(\ln2\gamma_K-\frac12\right)
\right].
\]

在固定 \(\omega,I_p\) 时，\(\gamma_K\propto E_0^{-1}\)。因此若把与 \(E_0\) 无关的场强尺度记为 \(E_*\)，主场强依赖可写成

\[
W_{\rm exp}\propto
\left(\frac{E_0}{E_*}\right)^{2I_p/(\hbar\omega)},
\]

其中 \(E_*\) 被吸收到指数精度之外的常数和预因子中。

如果最少需要吸收

\[
N_0=\left\lceil\frac{I_p+U_p}{\hbar\omega}\right\rceil
\]

个光子才能越过 dressing 后的阈值，那么微扰论给出更离散、更物理的通道标度

\[
M_{N_0}\propto E_0^{N_0},
\qquad
W_{N_0}\propto E_0^{2N_0}\propto\mathcal I^{N_0}.
\]

连续的 Keldysh 渐近指数中的 \(2I_p/(\hbar\omega)\) 因而应理解为与离散多光子阶数相接的半经典极限，而不是说实际光子数可以取任意实数。

\end{example}

\begin{note}[强场近似的适用条件]

\(\gamma_K\) 是最重要的第一控制参数，但不是完整相图。至少还要同时检查：

\[
\frac{U_p}{I_p},
\qquad
\frac{E_0}{E_{\rm BSI}},
\qquad
\frac{\hbar\omega}{I_p},
\qquad
ka_0,
\]

其中 \(E_{\rm BSI}\) 表示 势垒压低电离的特征场强。当外场已经把 Coulomb 势垒压到束缚能以下时，单纯谈“隧穿”就失去原来的禁阻区含义；当 \(ka_0\) 不再很小，偶极近似也必须修正；对短脉冲还需要把包络变化时间加入控制参数。因此 \(\gamma_K\ll1\) 并不自动保证所有准静态近似都成立。

\end{note}

\begin{example}[Coulomb 势垒的压低场强]
在静态、单电子 Coulomb 尾势模型中，沿电场降低势垒的方向取 $x>0$，写
\[
V(x)=-\frac{C_Z}{x}-eE_0x,
\qquad C_Z=\frac{Ze^2}{4\pi\epsilon_0},
\]
其中 $Z$ 为剩余离子的有效电荷数，束缚能为 $-I_p$。
势垒顶满足 $V'(x_s)=C_Z/x_s^2-eE_0=0$，所以
\[
x_s=\sqrt{\frac{C_Z}{eE_0}},\qquad
V(x_s)=-2\sqrt{C_ZeE_0}.
\]
当势垒顶降到束缚能时，静态禁阻区闭合，得到
\[
E_{\rm BSI}=\frac{I_p^2}{4C_Ze}
=\frac{\pi\epsilon_0 I_p^2}{Ze^3}.
\]
以氢基态 $Z=1$、$I_p=13.6\,\mathrm{eV}$ 代入，
$E_{\rm BSI}\simeq3.21\times10^{10}\,\mathrm{V/m}$。
若把 $E_0$ 视为正弦光场的峰值，则对应的周期平均强度为
\[
\mathcal I_{\rm BSI}=\frac12 c\epsilon_0 E_{\rm BSI}^2
\simeq1.37\times10^{14}\,\mathrm{W/cm^2}.
\]
这给出了 $E_0/E_{\rm BSI}$ 的实际尺度。多电子屏蔽、动态 Stark 位移和短脉冲
会改变阈值；该估计用于判断静态 Coulomb 模型中是否仍存在隧穿势垒，
不能替代含时电离率的计算。
\end{example}

\subsection{周期驱动的 Floquet--Sambe 结构}

若哈密顿量严格周期，

\[
\hat H(t+T)=\hat H(t),
\qquad T=\frac{2\pi}{\omega},
\]

则含时 Schr\"odinger 方程的形式精确结构由 Floquet 定理给出：

\begin{equation}
|\psi_\alpha(t)\rangle
=\ee^{-\ii\varepsilon_\alpha t/\hbar}
|u_\alpha(t)\rangle,
\qquad
|u_\alpha(t+T)\rangle=|u_\alpha(t)\rangle.
\label{eq:floquet-form}
\end{equation}

代回动力学方程得到扩展 Hilbert 空间中的本征方程

\begin{equation}
[\hat H(t)-\ii\hbar\partial_t]
|u_\alpha(t)\rangle
=\varepsilon_\alpha|u_\alpha(t)\rangle.
\label{eq:floquet-eigenproblem}
\end{equation}

准能量只在模 \(\hbar\omega\) 的意义下定义，因为

\[
|u_\alpha(t)\rangle
\longrightarrow
 \ee^{\ii n\omega t}|u_\alpha(t)\rangle
\]

会把

\[
\varepsilon_\alpha
\longrightarrow
\varepsilon_\alpha+n\hbar\omega
\]

而不改变物理态 \(|\psi_\alpha(t)\rangle\)。如果把周期函数进一步作 Fourier 展开

\[
|u_\alpha(t)\rangle
=\sum_{n\in\mathbb Z}
 \ee^{\ii n\omega t}|u_{\alpha n}\rangle,
\]

则 \eqnrefs{eq:floquet-eigenproblem} 变成“物理 Hilbert 空间 \(\otimes\) Fourier 格点”的无限维定态本征问题。这才是穿衣态与光子 sideband 的一般数学结构。

在强场相位中常遇到

\[
\exp[\ii z\sin\omega t].
\]

Jacobi--Anger 恒等式

\[
\exp(\ii z\sin\theta)
=
\sum_{n=-\infty}^{\infty}
J_n(z)\ee^{\ii n\theta}
\]

把它分解为整数倍频率通道。于是一个裸能级在 Floquet 表示中出现 \(E_0+n\hbar\omega\) 的重复副本。ATI 峰列可由这种准能量通道结构理解；近共振的 Autler--Townes 劈裂则对应在相关 Floquet 子空间中保留耦合并进行简并或近简并对角化。

\subsection{强场电离的实验参数与 Keldysh 区域}\label{sec:strong-field-experiments}

Keldysh 参数 \(\gamma_K=\sqrt{I_p/(2U_p)}\) 区分多光子（\(\gamma_K\gg1\)）与隧穿（\(\gamma_K\ll1\)）区域，其中 \(I_p\) 是电离势，\(U_p=e^2E_0^2/(4m\omega^2)\) 是有质动力能。

\begin{center}
\begin{tabular}{lcccc}
\hline
原子/分子 & \(I_p\) (eV) & 激光 & \(I\) (W/cm\(^2\)) & \(\gamma_K\) \\
\hline
H (1s) & 13.6 & 800\,nm & \(10^{14}\) & 1.07 \\
H (1s) & 13.6 & 800\,nm & \(10^{15}\) & 0.34 \\
He (1s\(^2\)) & 24.6 & 800\,nm & \(10^{15}\) & 0.45 \\
Ar (3p) & 15.8 & 800\,nm & \(10^{14}\) & 1.15 \\
Xe (5p) & 12.1 & 800\,nm & \(10^{14}\) & 1.01 \\
Xe (5p) & 12.1 & 800\,nm & \(10^{15}\) & 0.32 \\
H\(_2\) & 15.4 & 800\,nm & \(10^{14}\) & 1.14 \\
\hline
\end{tabular}
\end{center}

\begin{exbox}[Xe 的 ATI 谱]
Xe 在 800\,nm、\(1.4\times10^{14}\,\mathrm{W/cm^2}\)（\(\gamma_K\approx0.85\)）的典型 ATI 谱：光子能量 \(\hbar\omega=1.55\,\mathrm{eV}\)，有质动力能 \(U_p\approx8.4\,\mathrm{eV}\)。ATI 峰出现在 \(E_n=n\hbar\omega-I_p-U_p\)（阈上条件 \(n\hbar\omega>I_p+U_p\)，最小吸收光子数 \(n_{\min}=\lceil(I_p+U_p)/\hbar\omega\rceil=14\)），阈上第一峰 \(E_{14}\approx1.2\,\mathrm{eV}\)，随后 \(E_{15}\approx2.8\,\mathrm{eV}\)、\(E_{16}\approx4.3\,\mathrm{eV}\)。峰间距 \(\hbar\omega=1.55\,\mathrm{eV}\) 直接测量激光频率。高阶峰的强度受通道 closing 影响：当 \(U_p\) 增大到使 \(n\hbar\omega<I_p+U_p\) 时，该通道关闭，ATI 谱截断。
\end{exbox}

Keldysh 理论给出电离率的指数部分，隧穿区 $\gamma_K\ll1$ 的完整电离率由 ADK 公式给出，包含预因子和角分布。下面先把指数因子整理成标准形式并完整定义每个量，再说明预因子需要哪些额外输入。

定义三个基本量：电离势 $I_p>0$（把电子从基态移到连续阈所需的最小能量）、隧穿波数
\begin{equation}
\kappa\equiv\frac{\sqrt{2m_eI_p}}{\hbar},
\label{eq:adk-kappa}
\end{equation}
（量纲为 $1/$长度，是类氢束缚态尾部 $\psi_0(r)\sim r^{n^*-1}\ee^{-\kappa r}$ 的衰减常数），以及由束缚能设定的特征场强
\begin{equation}
E_c\equiv\frac{\hbar^2\kappa^3}{m_e e}.
\label{eq:adk-critical-field}
\end{equation}
量纲检验：$\hbar^2/m_e$ 是能量$\cdot$长度$^2$，乘 $\kappa^3$（长度$^{-3}$）得力，除以 $e$ 得场强。对氢 $1s$（$I_p=13.6\,\mathrm{eV}$），$\kappa=1/a_0$，$E_c=5.14\times10^{11}\,\mathrm{V/m}$，恰为原子单位场强。静场 WKB 指数\eqnrefs{eq:static-tunnel-exponent}——即禁阻区内 $\frac{2}{\hbar}\int_0^{I_p/(eE_0)}|p(x)|\dd x$，$|p(x)|=\sqrt{2m_e(I_p-eE_0x)}$——用这两个量可紧凑地写成
\[
\frac{2(2m_eI_p)^{3/2}}{3m_e e\hbar E_0}=\frac{2E_c}{3E_0}.
\]

在 \(\gamma_K\ll1\) 的隧穿区域，Ammosov--Delone--Krainov (ADK) 公式给出总电离率
\begin{equation}
W_{\rm ADK}=|C_{n^*l^*}|^2\,f(l,m)\,\frac{I_p}{\hbar}
\left(\frac{2E_c}{E_0}\right)^{2n^*-|m|-1}
\exp\!\left(-\frac{2E_c}{3E_0}\right),
\label{eq:adk-rate}
\end{equation}
其中各量的显式定义为
\[
n^*=\frac{Z}{\kappa a_0},
\qquad
l^*=n^*-1,
\qquad
|C_{n^*l^*}|^2=\frac{2^{2n^*}}{n^*\,\Gamma(n^*+l^*+1)\,\Gamma(n^*-l^*)},
\]
\[
f(l,m)=\frac{(2l+1)(l+|m|)!}{2^{|m|}\,|m|!\,(l-|m|)!}.
\]
这里 $Z$ 是电离后母离子的电荷数，$a_0$ 是 Bohr 半径，$n^*$ 是有效主量子数（$n^*=Z/\sqrt{2I_p}$ 取原子单位时的等价写法），$l,m$ 是初态轨道与磁量子数，$\Gamma$ 为 Euler Gamma 函数。$|C_{n^*l^*}|^2$ 是束缚态径向波函数渐近尾部 $R(r)\to C_{n^*l^*}\kappa^{3/2}(2\kappa r)^{n^*-1}\ee^{-\kappa r}$ 的归一化系数；$f(l,m)$ 来自角向波函数穿过势垒出口点时的权重。对氢 $1s$：$n^*=1$、$l=m=0$、$|C_{10}|^2=4$、$f(0,0)=1$，\eqnrefs{eq:adk-rate}退化为教科书中的氢原子静态隧穿率。

关于严格性的说明：指数因子 $\exp(-2E_c/3E_0)$ 只依赖禁阻区作用量，已由本章的 WKB/复时间鞍点分析完整导出。预因子则要求更多输入——必须把束缚态在核外的渐近形式与禁阻区内的 WKB 波函数在两者的交叠区匹配（这给出 $|C_{n^*l^*}|^2$ 与 $(2E_c/E_0)^{2n^*-|m|-1}$），并计算鞍点周围的二次涨落（这给出角分布因子 $f(l,m)$ 与速率量纲 $I_p/\hbar$）。这一渐近匹配由 Perelomov--Popov--Terent'ev (PPT) 理论完成、由 Ammosov--Delone--Krainov 化为上述显式形式；本书把它作为已证定理引用（推导见 Popov 的综述 Phys. Usp. \textbf{47}, 855 (2004)），但强调：指数部分是普遍且精确的，预因子依赖初态结构，两者必须分开检验。

\begin{exbox}[Ar 和 Xe 的隧穿电离率]
Ar (\(I_p=15.8\,\mathrm{eV}\)，\(n^*=\sqrt{13.6/15.8}=0.93\)) 在 \(10^{14}\,\mathrm{W/cm^2}\)、800\,nm：\(W_{\rm ADK}\approx3.1\times10^{13}\,\mathrm{s^{-1}}\)，对应寿命 \(\tau\approx32\,\mathrm{fs}\)。Xe (\(I_p=12.1\,\mathrm{eV}\)，\(n^*=\sqrt{13.6/12.1}=1.06\)) 在同场强：\(W_{\rm ADK}\approx8.5\times10^{14}\,\mathrm{s^{-1}}\)，\(\tau\approx1.2\,\mathrm{fs}\)。Xe 电离率远高于 Ar，解释了为什么 Xe 是高次谐波产生 (HHG) 的首选气体——足够快的电离提供足够的连续态电子，但又不至于在脉冲上升沿完全电离。
\end{exbox}

电离电子在激光场中返回核附近时可复合并发射高次谐波。三步模型（电离-加速-复合）给出截止频率的简单估计。

HHG 截止频率由三步模型的经典最大动能给出：
\[
\hbar\omega_{\rm cut}=I_p+3.17\,U_p
\]

\begin{center}
\begin{tabular}{lcccc}
\hline
气体 & \(I_p\) (eV) & \(\lambda\) (nm) & \(I\) (W/cm\(^2\)) & \(\hbar\omega_{\rm cut}\) (eV) \\
\hline
He & 24.6 & 800 & \(5\times10^{14}\) & 119.3 \\
Ne & 21.6 & 800 & \(5\times10^{14}\) & 116.3 \\
Ar & 15.8 & 800 & \(2\times10^{14}\) & 53.7 \\
Xe & 12.1 & 800 & \(1\times10^{14}\) & 31.0 \\
He & 24.6 & 1600 & \(2\times10^{14}\) & 176.1 \\
\hline
\end{tabular}
\end{center}

\begin{exbox}[He 在 1600\,nm 的水窗 HHG]
He 在 1600\,nm、\(2\times10^{14}\,\mathrm{W/cm^2}\) 的截止能量 \(\hbar\omega_{\rm cut}\approx176\,\mathrm{eV}\)，对应波长 \(\lambda\approx1240/176\approx7.0\,\mathrm{nm}\)，已接近水窗（280--540\,eV，即 2.3--4.4\,nm）的长波端。使用更短波长（如 2000\,nm 中红外）可进一步推高 \(U_p\propto\lambda^2\)，使截止进入水窗区域——这是中红外 HHG 用于生物成像的物理基础。
\end{exbox}

\begin{proof}[解]
第一步：从 SFA 波函数到含时偶极矩。取未耗尽近似下的 SFA 态
\begin{equation*}
|\psi(t)\rangle=\ee^{\ii I_pt/\hbar}|\psi_0\rangle
+\int\dd^3\vb p\;b_{\vb p}(t)\,|\chi_{\vb p}^V(t)\rangle,
\end{equation*}
其中第二项是电离产生的 Volkov 波包，$b_{\vb p}(t)$ 由直接电离振幅\eqnrefs{eq:sfa-direct-amplitude}的时间积分给出。偶极矩 $\vb d(t)=\langle\psi(t)|(-e\vb r)|\psi(t)\rangle$ 含三类项：基态静态偶极（不振荡，与谐波无关）、Volkov--Volkov 项（没有与基态的共振相位，只贡献缓变背景）、基态--连续态交叉项——只有它在基态能量与连续态能量的拍频上振荡，是 HHG 辐射的来源：
\begin{equation*}
\vb d(t)\simeq\int\dd^3\vb p\;\vb d_{\rm rec}^{*}(\vb p)\,b_{\vb p}(t)\,\ee^{-\ii I_pt/\hbar}+\mathrm{c.c.},
\qquad
\vb d_{\rm rec}(\vb p)\equiv\langle\chi_{\vb p}^V|(-e\vb r)|\psi_0\rangle.
\end{equation*}
代入 $b_{\vb p}$ 的积分表达式，把复合时刻记为 $t''$、电离时刻记为 $t'$，并作 Fourier 变换 $\vb d(\omega)=\int\vb d(t)\ee^{\ii\omega t}\dd t$，即得 Lewenstein 积分
\begin{equation*}
\vb d(\omega)=-\ii\int_0^t\dd t'\int_{t'}^t\dd t''\int\dd^3\vb p\;\vb d_{\rm rec}^*(\vb p+e\vb A(t''))\vb d_{\rm ion}(\vb p+e\vb A(t'))\ee^{-\ii S(\vb p,t',t'')/\hbar},
\end{equation*}
其中 $\vb d_{\rm ion}$ 是含场因子的电离矩阵元，半经典作用量
\begin{equation*}
S(\vb p,t',t'')=\int_{t'}^{t''}\left[\frac{(\vb p+e\vb A(\tau))^2}{2m_e}+I_p\right]\dd\tau-\hbar\omega t''
\end{equation*}
正是"电子在 $t'$ 进入连续态、在场中传播到 $t''$、复合时把能量差交给频率 $\omega$ 的光子"这条路径的相位。三个物理步骤各自对应积分的一个因子：$t'$ 处的电离矩阵元、$t'$ 到 $t''$ 的 Volkov 传播相位、$t''$ 处的复合矩阵元——这就是三步模型在量子理论中的精确含义，而不是外加的经典假设。

第二步：鞍点方程。谐波阶数高时相位 $S/\hbar$ 快速振荡，三重积分由驻相点主导。$\partial S/\partial\vb p=0$、$\partial S/\partial t'=0$、$\partial S/\partial t''=0$ 分别给出
\begin{align*}
\int_{t'_s}^{t''_s}\frac{\vb p_s+e\vb A(\tau)}{m_e}\,\dd\tau&=0
&&\text{(返回条件)},\\
\frac{[\vb p_s+e\vb A(t'_s)]^2}{2m_e}+I_p&=0
&&\text{(电离条件)},\\
\frac{[\vb p_s+e\vb A(t''_s)]^2}{2m_e}+I_p&=\hbar\omega
&&\text{(复合能量守恒)}.
\end{align*}
第一式要求电子在 $t''_s$ 时刻净位移为零、回到母离子；它同时固定中间正则动量 $\vb p_s=-\dfrac{e}{t''_s-t'_s}\displaystyle\int_{t'_s}^{t''_s}\vb A(\tau)\dd\tau$。第二式左端在实 $t'$ 上恒正，故 $t'_s$ 必为复数且 $\operatorname{Im}t'_s>0$，即隧穿电离的复时间，与单时间鞍点\eqnrefs{eq:sfa-saddle}同源。第三式是复合时刻把"动能$+$电离势"交给单个光子的能量守恒。

第三步：量子轨道与截止。在每个鞍点作二次展开并完成高斯积分，三重积分化为离散求和
\begin{equation*}
\vb d(\omega)\approx\sum_k C_k(\omega)\exp[-\ii S_k(\omega)/\hbar],
\end{equation*}
每组 $(t'_{s,k},t''_{s,k},\vb p_{s,k})$ 称为一条量子轨道：复时间隧穿、实轨道返回、单光子复合。把 $I_p\to0$ 并把鞍点取到实轴上，返回条件与能量守恒退化为\chapref{ch:strong-field-classical}的经典再碰撞方程；对单色场 $E_0\cos\omega t$ 在返回约束下最大化返回动能给出 $K_{\max}=3.17\,U_p$（电离相位 $\omega t'_s\approx17^\circ$、返回相位 $\omega t''_s\approx255^\circ$），于是量子理论给出同一条截止律
\begin{equation*}
\hbar\omega_{\rm cut}=I_p+3.17\,U_p.
\end{equation*}
多轨道相位差 $\Delta S=S_k-S_{k'}$ 产生谱内干涉（长、短轨道），可由双色场等脉冲整形控制。SFA 忽略的母态耗尽、Coulomb 传播与宏观相位匹配不改变鞍点结构，但定量预言强度时须另行加入。

例（HHG 截止的数值估计）。Ar 在 800\,nm、$10^{14}\,\mathrm{W/cm^2}$：用换算式 $U_p[\mathrm{eV}]=9.33\times10^{-14}\,I[\mathrm{W/cm^2}]\,\lambda^2[\mu\mathrm m^2]$ 得 $U_p=9.33\times10^{-14}\times10^{14}\times0.64=5.97\,\mathrm{eV}$，于是
\begin{equation*}
\hbar\omega_{\rm cut}=15.8+3.17\times5.97=34.7\,\mathrm{eV},
\qquad
N_{\rm cut}=\frac{34.7\,\mathrm{eV}}{1.55\,\mathrm{eV}}\approx22,
\end{equation*}
即约第 22 次谐波，与前表一致。光强提高到 $2\times10^{14}\,\mathrm{W/cm^2}$ 时 $U_p$ 加倍，截止升至 $53.7\,\mathrm{eV}$（第 35 次谐波）——截止随光强与 $\lambda^2$ 线性增长，是三步模型最直接的实验检验。

外场连续态基底由 Volkov 解给出。平面波场除 $\vb E$ 外还携带磁场 $\vb B=\hat{\vb e}_z(E_0/c)\sin\omega t$；电子的完整非相对论描述是 Pauli 哈密顿量，它在最小耦合之外还含自旋--磁场耦合 $-\vb\mu\cdot\vb B$（电子磁矩 $\vb\mu=-g\mu_B\vb S/\hbar$，$\mu_B=e\hbar/2m_e$ 为 Bohr 磁子，$g\approx2$），展开后即 $(e\hbar/2m_e)\vb\sigma\cdot\vb B$。保留该项时，平面波场中的非相对论 Volkov 解（原子单位）为
\begin{equation}
  \psi_{\boldsymbol p}(t,\boldsymbol x)=\exp\!\left[\ii\boldsymbol p\cdot\boldsymbol x-\ii\int_0^t\left(\frac{(\boldsymbol p+e\boldsymbol A(t'))^2}{2m}+\frac{e\,\boldsymbol\sigma\cdot\boldsymbol B(t')}{2m}\right)\dd t'\right].
  \label{eq:phys-sf-volkov-nonrel}
\end{equation}
自旋项的量级可由它与机械能项之比估计：$\dfrac{eB}{m}\Big/\dfrac{eAp}{m}\sim\dfrac{\omega}{pc}$，对光频场与非相对论光电子远小于 1；因此本章后续鞍点分析一律略去它，只在讨论自旋分辨电离（如 spin-polarized HHG）时才需恢复。SFA 即将波函数展开为束缚态加 Volkov 连续态。Coulomb 对连续态的修正在低能区显著，主要影响谱形而不改截止能 \(3.17U_p\)。
\end{proof}







\par
\begin{workedexample}[经典返回条件与 $3.17\,U_p$]
电子在 $t_i$ 时刻以零速度进入单色场 $\vb E(t)=\vb E_0\cos\omega t$。（i）积分 Newton 方程，证明返回母离子（$x(t_r)=0$）的条件是
\begin{equation*}
\cos\varphi_r-\cos\varphi_i+(\varphi_r-\varphi_i)\sin\varphi_i=0,
\qquad \varphi\equiv\omega t;
\end{equation*}
（ii）证明返回动能为 $K=2U_p(\sin\varphi_r-\sin\varphi_i)^2$；（iii）数值求解约束最大值，验证 $K_{\max}=3.17\,U_p$。

\textbf{解.} 积分 \(\ddot x=-(eE_0/m_e)\cos\omega t\)，并取 \(x(t_i)=\dot x(t_i)=0\)，得
\[
\dot x(t)=\frac{eE_0}{m_e\omega}(\sin\varphi_i-\sin\varphi),\quad
x(t)=\frac{eE_0}{m_e\omega^2}[\cos\varphi-\cos\varphi_i+(\varphi-\varphi_i)\sin\varphi_i].
\]
令 \(x(t_r)=0\) 即得题设返回条件。返回动能 \(m_e\dot x(t_r)^2/2=2U_p(\sin\varphi_r-\sin\varphi_i)^2\)。对返回条件隐式求导后令 \(\dd K/\dd\varphi_i=0\)，数值解为 \(\varphi_i\simeq0.31\)、\(\varphi_r\simeq4.40\)，代回得 \(K_{\max}\simeq3.17U_p\)。
\end{workedexample}
\par

\par
\begin{workedexample}[ADK 指数的场强敏感性]
对氢原子（$E_c=5.14\times10^{11}\,\mathrm{V/m}$），计算外场 $E_0=0.05E_c$ 与 $E_0=0.10E_c$ 时的 ADK 指数因子 $\exp(-2E_c/3E_0)$，并由此说明为什么强场电离率对光强的依赖几乎是"开关式"的。

\textbf{解.} 两个指数分别为
\[
\exp\!\left(-\frac{2E_c}{3(0.05E_c)}\right)=\ee^{-13.333}\simeq1.62\times10^{-6},\qquad
\exp\!\left(-\frac{2E_c}{3(0.10E_c)}\right)=\ee^{-6.667}\simeq1.27\times10^{-3}.
\]
场强加倍使指数因子增大约 \(7.8\times10^2\) 倍，尚未计入前置幂律因子；这就是电离率对场强呈强烈非线性响应的原因。
\end{workedexample}
\par
